\section{Introducción}

La IA aplicada al agro incluye algoritmos de machine learning para analizar datos agrícolas, predecir patrones y optimizar recursos. En Colombia, la IA puede revolucionar prácticas tradicionales, reduciendo pérdidas y aumentando eficiencia.

El Portal Agro Comercial del Huila incorpora IA en módulos como diagnóstico de plagas y recomendaciones de siembra, haciendo estas tecnologías accesibles a todos los productores.

\section{Marco teórico y trabajos relacionados}

Literatura sobre IA en agricultura incluye estudios de la FAO (2023) sobre predicción de cosechas y análisis de imágenes (2024). Investigaciones en "AI for Sustainable Agriculture" (2025) destacan beneficios en precisión y sostenibilidad.

Trabajos relacionados cubren ética de IA en agro (2024) y casos en América Latina (CEPAL, 2020).

\section{Metodología}

\subsection{Desarrollo de Modelos IA}
Se entrenaron modelos con datos locales del Huila para predicciones precisas.

\subsection{Integración en Plataforma}
Módulos IA integrados en el portal, con interfaces amigables.

\subsection{Validación}
Pruebas con productores para medir precisión y usabilidad.

\section{Implementación del Portal Agro Comercial del Huila}

El portal usa IA para análisis de datos, pronósticos y asesoría.

\subsection{Aplicaciones Específicas}
Diagnóstico de plagas, optimización de riego, predicción de precios.

\subsection{Acceso Inclusivo}
Herramientas gratuitas y capacitación para uso efectivo.

\section{Resultados e Impactos}

La IA ha mejorado la precisión en diagnósticos en un 80\%, reduciendo pérdidas.

\subsection{Eficiencia en Toma de Decisiones}
Productores toman decisiones basadas en datos, aumentando rendimientos.

\subsection{Adopción}
Alta aceptación entre usuarios capacitados.

\section{Discusión}

\subsection{Desafíos}
Acceso a datos y sesgos en modelos.

\subsection{Comparación con Literatura}
Resultados alineados, confirmando potencial de IA.

\section{Conclusiones y Visión Futura}

La IA es un aliado clave en el agro moderno. El portal facilita su adopción, impulsando innovación.

\subsection{Trabajo Futuro}
Mejorar modelos y expandir aplicaciones.

\section{Implementación en el proyecto Portal Agrocomercial del Huila}

El Portal Agrocomercial del Huila integra inteligencia artificial para proporcionar diagnósticos precisos de plagas mediante análisis de imágenes, predicciones de cosechas basadas en datos históricos y recomendaciones personalizadas de manejo agrícola. Estas herramientas permiten a los productores optimizar recursos y reducir riesgos, democratizando el acceso a tecnologías avanzadas. Además, el portal incluye capacitaciones en el uso de IA, asegurando que los beneficios lleguen a comunidades rurales sin barreras técnicas.